% !TEX TS-program = pdflatex
% !TeX program = pdflatex
% !TEX encoding = UTF-8
% !TeX spellcheck = en_US

\documentclass[xcolor=table]{beamer}


%\usepackage{fullpage}
%\usepackage[left=2.8cm,right=2.2cm,top=2 cm,bottom=2 cm]{geometry}
\setbeamersize{text margin left=10pt,text margin right=10pt}
\usepackage{amsmath,amssymb} 
\usepackage{fontenc}
\usepackage{textcomp}
\usepackage[utf8]{inputenc}
\usepackage[frenchb]{babel}
\usepackage{arabtex,acolor}
\usepackage{txfonts}
\usepackage{tipa}
\usepackage{CJKutf8}
\usepackage[]{graphicx}
\usepackage{multirow}
\usepackage{hyperref}
\usepackage{colortbl}
\usepackage{tabularray}
\usepackage{listingsutf8}
\usepackage{wrapfig}
\usepackage{multicol}
\usepackage[export]{adjustbox} %for images in table, also for frame
\usepackage[many]{tcolorbox}
\usepackage{wasysym}
\usepackage[french,lined]{algorithm2e}
\usepackage{alltt} %verbatim with commands
\usepackage{amsmath}

\tcbuselibrary{listings}

%\hypersetup{
%	colorlinks,
%	urlcolor = blue
%}

%\renewcommand{\baselinestretch}{1.5}

\def\supit#1{\raisebox{0.8ex}{\small\it #1}\hspace{0.05em}}

\AtBeginSection{%
	\begin{frame}
		\sectionpage
	\end{frame}
}

\newcommand{\rottext}[2]{%
	\rotatebox{90}{%
	\begin{minipage}{#1}%
		\raggedleft#2%
	\end{minipage}%
	}%
}

\usepackage{longtable}
\usepackage{tabu}


\institute{ %
	\parbox{.3\textwidth}{\begin{center}
			Laboratoire de la Communication dans les Systèmes Informatiques (LCSI)\\
			\includegraphics[height=.8cm]{../extra/logo/lcsi-logo.png}
	\end{center}}
	\hskip2cm%\hfill
	\parbox{.3\textwidth}{\begin{center}
			Laboratoire des Méthodes de Conception des Systèmes	(LMCS)\\
			\includegraphics[height=1cm]{../extra/logo/lmcs-logo.png}
		\end{center}
	}
	
	
	\vspace{6pt}
	École nationale Supérieure d'Informatique (ESI, ex. INI), Alger, Algérie
}

\author[ESI - ARIES \& BENATCHBA (2022-2023)] %
{Abdelkrime Aries \hskip2cm Karima Benatchba}

\titlegraphic{%
	\includegraphics[height=1cm]{../extra/logo/esi-logo.png}%
	\hspace*{2cm}%
	\includegraphics[height=1cm]{../extra/logo/esi.ml.pdf}%
}%\hspace*{4.75cm}~

\date{Academic year: 2022/2023} %\today

\usetheme{Karim} % Antibes Boadilla Warsaw

\beamertemplatenavigationsymbolsempty

\let\oldcite\cite
\renewcommand{\cite}[1]{{\bfseries\color{red}\oldcite{#1}}}

\let\oldurl\url
\renewcommand{\url}[1]{{\color{blue}\oldurl{#1}}}

%\setbeamertemplate{headline}{}

\definecolor{lightblue}{HTML}{D0D2FF}
\definecolor{lightyellow}{HTML}{FFFFAA}
\definecolor{darkblue}{HTML}{0000BB}
\definecolor{olivegreen}{HTML}{006600}
\definecolor{violet}{HTML}{6600CC}

\newcommand{\keyword}[1]{\textcolor{red}{\bfseries\itshape #1}}
\newcommand{\expword}[1]{\textcolor{olivegreen}{#1}}
\newcommand{\optword}[1]{\textcolor{violet}{\bfseries #1}}

\makeatletter
\newcommand\mysphere{%
	\parbox[t]{10pt}{\raisebox{0.2pt}{\beamer@usesphere{item projected}{bigsphere}}}}
\makeatother

%\let\oldtabular\tabular
%\let\endoldtabular\endtabular
%\renewenvironment{tabular}{\rowcolors{2}{white}{lightblue}\oldtabular\rowcolor{blue}}{\endoldtabular}


\NoAutoSpacing %french autospacing after ":"

\def\graphpath{}

\newcommand{\changegraphpath}[1]{\def\graphpath{#1}}


\newcommand{\vgraphpage}[2][.84\textheight]{%
%	\begin{center}%
		\includegraphics[height=#1]{\graphpath #2}%
%	\end{center}%
}

\newcommand{\hgraphpage}[2][\textwidth]{%
%	\begin{center}%
		\includegraphics[width=#1]{\graphpath #2}%
%	\end{center}%
}

\newcommand{\graphpage}[2][]{%
	\includegraphics[#1]{\graphpath #2}%
}

\bibliographystyle{apalike}

\newcommand{\insertbibliography}[2]{
	\appendix
	\section*{Bibliography}
	\nocite{#2}
%	\makeatletter % to change template
%	\setbeamertemplate{headline}[default] % not mandatory, but I though it was better to set it blank
%	\def\beamer@entrycode{\vspace*{-\headheight}} % here is the part we are interested in :)
%	\makeatother
	\begin{multicols*}{2}[\frametitle{\insertsection} \usebeamertemplate*{frametitle}]%\usebeamertemplate*{frametitle}\frametitle{Références}
		\tiny
		\bibliography{#1}
	\end{multicols*}
}

\definecolor{my-grey}{RGB}{233, 233, 233}

\newcommand{\insertlicence}{
	\begin{frame}[plain]
		\frametitle{License: CC-BY 4.0}
		
		\begin{tcolorbox}[colback=cyan,
			colframe=cyan,  
			arc=0pt,outer arc=0pt,
			valign=top, 
			halign=center,
			width=\textwidth]
			
			\includegraphics[width=.5cm]{../img/licence/cc_icon_white_x2.png}
			\includegraphics[width=.5cm]{../img/licence/attribution_icon_white_x2.png}
			
			\color{white}
			\bfseries Attribution 4.0 International (CC BY 4.0) \\
			\tiny \url{https://creativecommons.org/licenses/by/4.0/deed.en}
			
		\end{tcolorbox}\vspace{-.5cm}
		\begin{tcolorbox}[colback=my-grey,
			colframe=my-grey,  
			center, arc=0pt,outer arc=0pt,
			valign=top, 
			halign=left,
			width=\textwidth]
			
			\tiny
			
			\begin{center}
				\bfseries\Large
				You are free to:
			\end{center}
			
			\begin{minipage}{0.83\textwidth}
				\begin{itemize}
					\item[] \textbf{Share} — copy and redistribute the material in any medium or format
					\item[] \textbf{Adapt} — remix, transform, and build upon the material
					for any purpose, even commercially.
				\end{itemize}
				
				\vspace{12pt}
				
				\textit{The licensor cannot revoke these freedoms as long as you follow the license terms.}
			\end{minipage}
			\begin{minipage}{0.15\textwidth}
				\includegraphics[width=\textwidth]{../img/licence/FreeCulturalWorks_seal_x2.jpg}
			\end{minipage}
			
			
			\begin{center}
				\bfseries\Large
				Under the following terms:
			\end{center}
			
			\begin{itemize}
				\item[] \textbf{Attribution} — You must give appropriate credit, provide a link to the license, and indicate if changes were made. You may do so in any reasonable manner, but not in any way that suggests the licensor endorses you or your use.
				\item[] \textbf{No additional restrictions} — You may not apply legal terms or technological measures that legally restrict others from doing anything the license permits.
			\end{itemize}
			
		\end{tcolorbox}
	\end{frame}
}

\settowidth{\leftmargini}{\usebeamertemplate{itemize item}}
\addtolength{\leftmargini}{\labelsep}

\AtBeginDocument{
	\newcolumntype{L}[2]{>{\vbox to #2\bgroup\vfill\flushleft}p{#1}<{\egroup}} 
	
	\begin{frame}[plain]
		\maketitle
	\end{frame}
	
	\insertlicence
}


% needs etoolbox; to break links after -
\appto\UrlBreaks{\do\-}


\makeatletter
\newcommand{\xRightarrow}[2][]{\ext@arrow 0359\Rightarrowfill@{#1}{#2}}
\makeatother

\usefonttheme{structurebold}

% figure caption
\setlength\abovecaptionskip{2pt}

% equations
\AtBeginDocument{
	\abovedisplayskip=0pt
	\abovedisplayshortskip=0pt
	\belowdisplayskip=0pt
	\belowdisplayshortskip=0pt
}

\title[ML-TP: Multi-class classification] %
{Machine Learning (Lab support)\\Multi-class classification} 

\changegraphpath{../img/multi-class/}

\begin{document}
	
\begin{frame}
	\frametitle{Machine Learning}
	\framesubtitle{Multi-class classification: Introduction}
	
	\begin{itemize}
		\item We have already seen ...
		\begin{itemize}
			\item how to estimate the probability of a class
			\item using logistic regression
			\item in case of binary classification (belongs to a class or not)
		\end{itemize}
		\item But ...
		\begin{itemize}
			\item how to estimate them in case of multiple classes?
			\item how to affect a sample into many classes at once?
		\end{itemize}

	\end{itemize}
\end{frame}


\begin{frame}
	\frametitle{Machine Learning}
	\framesubtitle{Multi-class classification: Plan}
	
	\begin{multicols}{2}
		%	\small
		\tableofcontents
	\end{multicols}
\end{frame}

\section{Classification}

\begin{frame}
	\frametitle{Machine Learning}
	\framesubtitle{Classification}
	
	\hgraphpage{classif.pdf}
	
\end{frame}

\subsection{Binary classification}


\begin{frame}
	\frametitle{Machine Learning: Classification}
	\framesubtitle{Binary classification}
	
	\begin{minipage}{0.65\textwidth}
	\begin{itemize}
		\item Number of available classes: 2 (Actually, it is just one class. We test if a sample belongs to it or not)
		\item Number of chosen classes: 1
		\item Examples
		\begin{itemize}
			\item \expword{Classify an image as containing a specific object or not.}
			\item \expword{Finding out if a message is a spam or not.}
		\end{itemize}
	\end{itemize}
	\end{minipage}
	\begin{minipage}{0.3\textwidth}
		\hgraphpage{binary.pdf}
	\end{minipage}
	
\end{frame}

\subsection{Multi-class classification}

\begin{frame}
	\frametitle{Machine Learning: Classification}
	\framesubtitle{Multi-class classification}
	
	\begin{minipage}{0.65\textwidth}
	\begin{itemize}
		\item Number of available classes: L $ \ge $ 2
		\item Number of chosen classes: 1
		\item Examples
		\begin{itemize}
			\item \expword{Detect an animal from an input image.}
			\item \expword{Detect a movie's certification (PG, R, etc.) from its description.}
		\end{itemize}
	\end{itemize}
	\end{minipage}
	\begin{minipage}{0.3\textwidth}
		\hgraphpage{multiclass.pdf}
	\end{minipage}
	
\end{frame}

\begin{frame}
	\frametitle{Machine Learning: Classification}
	\framesubtitle{Multi-class classification: Methods}
	
	\hgraphpage{multi-class.pdf}
	
\end{frame}

\subsection{Multi-label classification}

\begin{frame}
	\frametitle{Machine Learning: Classification}
	\framesubtitle{Multi-label classification}
	
	\begin{minipage}{0.65\textwidth}
	\begin{itemize}
		\item Number of available classes: L $ \ge $ 2
		\item Number of chosen classes: K $ \le $ L
		\item Examples
		\begin{itemize}
			\item \expword{Detect many animals from an input image.}
			\item \expword{Find the clothes (hat, jeans, scarf, etc.) from an image.}
			\item \expword{Detect movie's genres (Sci-fi, Action, Comedy, etc.) from its description.}
		\end{itemize}
	\end{itemize}
	\end{minipage}
	\begin{minipage}{0.3\textwidth}
		\hgraphpage{multilabel.pdf}
	\end{minipage}
	
\end{frame}

\begin{frame}
	\frametitle{Machine Learning: Classification}
	\framesubtitle{Multi-label classification: Methods \cite{2012-madjarov-al}}
	
	\hgraphpage{multi-label.pdf}
	
	\begin{itemize}
		\item scikit-multilearn: \url{http://scikit.ml/}
	\end{itemize}
	
\end{frame}

\section{Binary logistic regression}

\begin{frame}
	\frametitle{Machine Learning}
	\framesubtitle{Binary logistic regression}
	
	\begin{center}
		\vgraphpage{LR_bin.pdf}
	\end{center}
	
\end{frame}

\subsection{Probability estimation}

\begin{frame}
\frametitle{Machine Learning: Binary logistic regression}
\framesubtitle{Probability estimation}

\[Z = \sum_{j=1}^{N} \theta_j X_j = X \cdot \theta\]
\begin{itemize}
	\item $Z[M]$ is a vector of $M$ elements (samples)
	\item $X[M, N]$ is a matrix of $M$ samples and $N$ features
	\item $\theta[N]$ is a vector of $N$ parameters (weights)
\end{itemize}

\[H = \sigma(Z) = \frac{1}{1+e^{-Z}}\]

\begin{itemize}
	\item $H[M]$ is a vector of $M$ probabilities (samples)
\end{itemize}

\end{frame}

\subsection{Cost and Gradient}

\begin{frame}
	\frametitle{Machine Learning: Binary logistic regression}
	\framesubtitle{Cost and Gradient}

	\[J_\theta = BCE = \frac{-1}{M} \sum\limits_{i=1}^{M} [Y^{(i)} \log(H^{(i)}) + (1- Y^{(i)}) \log(1 - H^{(i)})]\]
	\begin{itemize}
		\item $Y[M]$ and $H[M]$ are two vectors of $M$ elements (samples)
		\item $J_\theta$ is a scalar
	\end{itemize}
	
	
	\[
	\frac{\partial BCE}{\partial \theta_j} = \frac{1}{M} \sum\limits_{i=1}^{M} (H^{(i)} - Y^{(i)}) X_j^{(i)}
	\]
	
	\[
	\frac{\partial BCE}{\partial \theta} = \frac{1}{M} (H - Y) \cdot X
	\]
	
	\begin{itemize}
		\item $\frac{\partial BCE}{\partial \theta}[N]$ is a vector of $N$ elements (features)
	\end{itemize}
	
\end{frame}

\subsection{Gradient (derivation)}

\begin{frame}
	\frametitle{Machine Learning: Binary logistic regression}
	\framesubtitle{Gradient (derivation)}

	\tiny\vspace{-6pt}
	\begin{align*}
	\frac{\partial BCE}{\partial \theta_j} 
	& = & \frac{-1}{M} \sum\limits_{i=1}^{M} \frac{\partial}{\partial \theta_j} [Y^{(i)} \log(H^{(i)}) + (1- Y^{(i)}) \log(1 - H^{(i)})] \\
	& = & \frac{-1}{M} \sum\limits_{i=1}^{M} [ Y^{(i)} \frac{\partial}{\partial \theta_j} \log(H^{(i)}) + (1- Y^{(i)}) \frac{\partial}{\partial \theta_j}\log(1 - H^{(i)})]\\
	& = & \frac{-1}{M} \sum\limits_{i=1}^{M} [ Y^{(i)} \frac{1}{H^{(i)}} \frac{\partial}{\partial \theta_j} H^{(i)} + (1- Y^{(i)}) \frac{-1}{1-H^{(i)}} \frac{\partial}{\partial \theta_j} H^{(i)})] \\
	& = & \frac{-1}{M} \sum\limits_{i=1}^{M} \frac{Y^{(i)}-H^{(i)}}{H^{(i)}(1-H^{(i)})} \frac{\partial}{\partial \theta_j} H^{(i)} \\
	\end{align*}
	
	\vspace{-3pt}
	\begin{align*}
	\frac{\partial H^{(i)}}{\partial \theta_j} 
	& = & \frac{\partial \sigma(Z^{(i)})}{\partial Z^{(i)}} \frac{\partial Z^{(i)}}{\partial \theta_j}
	 =  [\sigma(Z^{(i)}) (1-\sigma(Z^{(i)}))]\frac{\partial}{\partial \theta_j} \sum\limits_{k=0}^{N} \theta_k X_k^{(i)} 
	 =  H^{(i)} (1-H^{(i)})  X_j^{(i)}\\
	\end{align*}
	
	\vspace{-3pt}
	\begin{align*}
	\frac{\partial BCE}{\partial \theta_j} 
	& = & \frac{-1}{M} \sum\limits_{i=1}^{M} \frac{Y^{(i)}-H^{(i)}}{H^{(i)}(1-H^{(i)})} [H^{(i)} (1-H^{(i)}) X_j^{(i)}] \\
	& = & \frac{1}{M} \sum\limits_{i=1}^{M} (H^{(i)} - Y^{(i)}) X_j^{(i)}\\
	\end{align*}
		
\end{frame}

\subsection{Parameters' update}

\begin{frame}
	\frametitle{Machine Learning: Binary logistic regression}
	\framesubtitle{Parameters' update}

	\[\theta = \theta - \alpha \frac{\partial J_\theta}{\partial \theta}\]
	
	\begin{itemize}
		\item $\frac{\partial J_\theta}{\partial \theta}[N]$ is a vector of $N$ elements (features)
		\item $\theta[N]$ is a vector of $N$ elements (features)
		\item $\alpha $ is a learning rate
	\end{itemize}
	
\end{frame}

\section{Multi-class logistic regression}

%\begin{frame}
%	\frametitle{Machine Learning}
%	\framesubtitle{Multi-class logistic regression}
%	
%	\hgraphpage{multi-class.pdf}
%	
%\end{frame}

\subsection{One-vs-Rest}

\begin{frame}
	\frametitle{Machine Learning: Multi-class logistic regression}
	\framesubtitle{One-vs-Rest}
	
	\begin{center}
		\vgraphpage{LR_ovr.pdf}
	\end{center}
	
\end{frame}

\begin{frame}
	\frametitle{Machine Learning: Multi-class logistic regression}
	\framesubtitle{One-vs-Rest: Description}
	
	Given $L$ output classes:
	\begin{itemize}
		\item Training
		\begin{itemize}
			\item For each class $C_l$, we train a binary model $M_l$ separately
			\item In this case, we will have $L$ binary models
			\item The positive class is represented by $C_l$ class's samples
			\item The negative class is represented by the rest of the samples
		\end{itemize}
		\item Estimation
		\begin{itemize}
			\item Given a sample
			\item For each class $C_l$, we estimate its probability using $M_l$
			\item We take the class with the maximum probability
			\item In this case, the sum of probabilities does not always give $1$
			\[\sum_{l=1}^{L} P(l|x; \theta^{M_l}) \in [0, L]\]
		\end{itemize}
	\end{itemize}
	
\end{frame}

\subsection{One-vs-One}

\begin{frame}
	\frametitle{Machine Learning: Multi-class logistic regression}
	\framesubtitle{One-vs-One}
	
%	\begin{center}
		\hgraphpage{LR_ovo.pdf}
%	\end{center}
	
\end{frame}

\begin{frame}
	\frametitle{Machine Learning: Multi-class logistic regression}
	\framesubtitle{One-vs-One: Description}
	
	Given $L$ output classes:
	\begin{itemize}
		\item Training
		\begin{itemize}
			\item For each two classes $C_{l}$ and $C_{l'}$, we train a binary model $M_{ll'}$ separately
			\item In this case, we will have $L (L-1)/2$ binary models
			\item The positive class is represented one of the two classes
			\item The negative class is represented by the other class
		\end{itemize}
		\item Estimation
		\begin{itemize}
			\item Given a sample
			\item For each model $C_{ll'}$, we estimate a probability
			\item In this case, the class is either $l$ or $l'$
			\item We count the number of each class being estimated
			\item The majority wins (vote)
		\end{itemize}
	\end{itemize}
	
\end{frame}

\subsection{Multinomial}

\begin{frame}
	\frametitle{Machine Learning: Multi-class logistic regression}
	\framesubtitle{Multinomial}
	
	\begin{center}
		\vgraphpage{LR_multi.pdf}
	\end{center}
	
\end{frame}

\begin{frame}
	\frametitle{Machine Learning: Multi-class logistic regression}
	\framesubtitle{Multinomial: Description}
	
	Given $L$ output classes:
	\begin{itemize}
		\item This is a generalization of binary classification in logistic regression
		\item Training
		\begin{itemize}
			\item For each class $C_{l}$, we calculate a weighted sum
			\item We apply \textbf{Softmax} function on these outputs
			\item It is a function which transforms the sums into probabilities
			\item Also, it normalizes these probabilities to get a sum of $1$
			\item The output classes are represented using One-Hot
			\item The cost function is a generalized version of that of binary classification
		\end{itemize}
		\item Estimation
		\begin{itemize}
			\item Given a sample
			\item We apply the model to get a vector of probabilities
			\item The most probable class wins
		\end{itemize}
	\end{itemize}
	
\end{frame}

\begin{frame}
	\frametitle{Machine Learning: Multi-class logistic regression}
	\framesubtitle{Multinomial: Probability estimation}
	
	\[Z = \sum_{j=1}^{N} \theta_j X_j = X \cdot \theta\]
	\begin{itemize}
		\item $Z[M, L]$ is a matrix of $M$ samples and $L$ classes
		\item $X[M, N]$ is a matrix of $M$ samples and $N$ features
		\item $\theta[N, L]$ is a matrix of $N$ features and $L$ classes
	\end{itemize}
	
	\[H = softmax(Z) = \frac{e^{Z}}{\sum_{k=1}^{L} e^{Z_k}}\]
	
	\begin{itemize}
		\item $H[M, L]$ is a matrix of $M$ samples and $L$ classes
	\end{itemize}
	
\end{frame}

\begin{frame}
	\frametitle{Machine Learning: Multi-class logistic regression}
	\framesubtitle{Multinomial: Cost and gradient}
	
	\[J_\theta = \frac{-1}{M} \sum\limits_{i=1}^{M} \sum_{k=1}^{L} Y^{(i)}_k \log(H^{(i)}_k)\]
	\begin{itemize}
		\item $Y[M, L]$ and $H[M, L]$ are two matrices of $M\times L$ elements (features X classes)
		\item $J_\theta$ is a scalar
	\end{itemize}
	
	
	\[
	\frac{\partial BCE}{\partial \theta_j} = \frac{1}{M} \sum\limits_{i=1}^{M} (H^{(i)} - Y^{(i)}) X_j^{(i)}
	\]
	
	\[
	\frac{\partial BCE}{\partial \theta} = \frac{1}{M} (H - Y) \cdot X
	\]
	
	\begin{itemize}
		\item $\frac{\partial BCE}{\partial \theta}[N, L]$ is a matrix of $N\times L$ elements (features X classes)
	\end{itemize}
	
\end{frame}

\begin{frame}
	\frametitle{Machine Learning: Multi-class logistic regression}
	\framesubtitle{Multinomial: Parameters' update}
	
	\[\theta = \theta - \alpha \frac{\partial J_\theta}{\partial \theta}\]
	
	\begin{itemize}
		\item $\frac{\partial BCE}{\partial \theta}[N, L]$ is a matrix of $N\times L$ elements (features X classes)
		\item $\theta[N, L]$ is a matrix of $N$ features and $L$ classes
		\item $\alpha$ is the learning rate
	\end{itemize}
	
\end{frame}

\section{Multi-label logistic regression}

\subsection{Binary relevance}

\begin{frame}
	\frametitle{Machine Learning: Multi-label logistic regression}
	\framesubtitle{Binary relevance}
	
	\begin{center}
		\vgraphpage{LR_binary_relevance.pdf}
	\end{center}
	
\end{frame}

\begin{frame}
	\frametitle{Machine Learning: Multi-class logistic regression}
	\framesubtitle{Binary relevance: Description}
	
	Given $L$ output classes:
	\begin{itemize}
		\item Training
		\begin{itemize}
			\item For each class $C_l$, we train a binary model $M_l$ separately
			\item In this case, we will have $L$ binary models
			\item The positive class is represented by $C_l$ class's samples
			\item The negative class is represented by the rest of the samples
		\end{itemize}
		\item Estimation
		\begin{itemize}
			\item Given a sample
			\item For each class $C_l$, we estimate its probability using $M_l$
			\item If the probability is greater or equals 50\%, then the sample belongs to the class $C_l$
		\end{itemize}
	\end{itemize}
	
\end{frame}

\subsection{Label powerset}

\begin{frame}
	\frametitle{Machine Learning: Multi-label logistic regression}
	\framesubtitle{Label powerset}
	
%	\begin{center}
		\hgraphpage{LR_label_powerset.pdf}
%	\end{center}
	
\end{frame}

\begin{frame}
	\frametitle{Machine Learning: Multi-class logistic regression}
	\framesubtitle{Label powerset: Description}
	
	Given $L$ output classes:
	\begin{itemize}
		\item Training
		\begin{itemize}
			\item We look for all combinations of classes
			\item Each combination is considered as a new class
			\item We train a multi-class model
		\end{itemize}
		\item Estimation
		\begin{itemize}
			\item Given a sample
			\item We use our trained multi-class model to get one class
			\item This class is a combination of many original classes
		\end{itemize}
	\end{itemize}
	
\end{frame}

\insertbibliography{EN.ML_Multi-class}{*}

\end{document}



